%%=============================================================================
%% Conclusie
%%=============================================================================

\chapter{Conclusie}%
\label{ch:conclusie}


Het doel van deze bachelorproef was het overbruggen van de kloof tussen fysieke medische simulatiehardware en 
immersieve softwaretechnologie, om zo tot een intuïtievere en effectievere reanimatietraining te komen. 
Binnen de huidige gezondheidszorgopleidingen is de kwaliteit van \textit{Cardiopulmonaire Resuscitatie} (CPR) 
direct gekoppeld aan overlevingskansen in reële noodsituaties. Hoewel geavanceerde oefenpoppen over nauwkeurige 
sensoren beschikken, blijft de feedback vaak beperkt tot abstracte interfaces op externe schermen.

Met de realisatie van dit onderzoek is een werkend \textit{Proof of Concept} (PoC) opgeleverd voor de Meta Quest 3. 
Door middel van een \textit{reverse engineering}-traject is het gesloten Bluetooth-protocol van een commerciële 
reanimatiepop succesvol ontsloten. Deze realtime datastroom is vervolgens geïntegreerd in een Mixed Reality-omgeving 
die de student via multisensorische feedback bestaande uit procedurele visuele hulpmiddelen, adaptieve audio en 
gamificatie-elementen direct op het fysieke interactiepunt stuurt. In dit afsluitende hoofdstuk worden de 
resultaten geëvalueerd aan de hand van de initiële deel- en hoofdonderzoeksvragen, gevolgd door een kwalitatieve 
expert-validatie en een kritische blik op de toekomstige ontwikkelingsmogelijkheden van het systeem.    TESTTEST

